\documentclass{article}
\usepackage{amsmath,amssymb,amsthm}
\usepackage{algorithm}
\usepackage{algpseudocode}
\usepackage{fullpage}
\usepackage{enumitem}
\usepackage{hyperref}
\newtheorem{theorem}{Theorem}
\newtheorem{lemma}{Lemma}
\newtheorem{assumption}{Assumption}
\newcommand{\sign}{\operatorname{sign}}
\newcommand{\E}{\mathbb{E}}
\newcommand{\R}{\mathbb{R}}

\begin{document}

\section*{SignSGD with Median‑Based Variance Reduction}

\section*{Mathematical Formulation}
Let $f:\R^{d}\rightarrow\R$ be differentiable.  
At iteration $k$ we draw a minibatch $\mathcal{B}_{k}$ of size $B$ and compute stochastic gradients
\[
g_{k}^{(i)}\;=\;\nabla f(x_{k};\xi_{k}^{(i)})\,,\qquad i=1,\dots ,B,
\]
where $\xi_{k}^{(i)}$ denotes the random data sample.  
Define the coordinate‑wise median of the signs
\[
m_{k}\;:=\;\operatorname{median}\bigl\{\sign(g_{k}^{(1)}),\dots ,\sign(g_{k}^{(B)})\bigr\}\in\{-1,0,1\}^{d}.
\]
The update rule of **SignSGD‑Median** is
\[
\boxed{x_{k+1}=x_{k}-\alpha\, m_{k}},\qquad k=0,1,\dots ,
\]
where $\alpha>0$ is a (possibly diminishing) stepsize.

\section*{Pseudocode}
\begin{algorithm}[H]
\caption{SignSGD‑Median}
\begin{algorithmic}[1]
\State \textbf{Input:} stepsize $\alpha$, minibatch size $B$, max. iter $T$
\State Initialize $x_{0}\in\R^{d}$
\For{$k=0$ \textbf{to} $T-1$}
    \State Sample a minibatch $\mathcal{B}_{k}=\{\xi_{k}^{(i)}\}_{i=1}^{B}$
    \State Compute stochastic gradients $g_{k}^{(i)}=\nabla f(x_{k};\xi_{k}^{(i)})$
    \State $m_{k}\gets\operatorname{median}\{\sign(g_{k}^{(1)}),\dots ,\sign(g_{k}^{(B)})\}$
    \State $x_{k+1}\gets x_{k}-\alpha\, m_{k}$
\EndFor
\State \textbf{Output:} $x_{T}$
\end{algorithmic}
\end{algorithm}

\section*{Dry Run Example}
Consider the one‑dimensional quadratic $f(w)=(w-3)^{2}$, $w_{0}=0$, stepsize $\alpha=0.1$, and a deterministic “minibatch’’ (so $m_{k}=\sign(\nabla f(w_{k}))$).

\begin{tabular}{c|c|c|c}
\textbf{Iter.} $k$ & $w_{k}$ & $\nabla f(w_{k})=2(w_{k}-3)$ & $w_{k+1}=w_{k}-0.1\sign(\nabla f(w_{k}))$ \\ \hline
0 & $0$ & $-6$ & $0-0.1\cdot(-1)=0.1$ \\
1 & $0.1$ & $-5.8$ & $0.1-0.1\cdot(-1)=0.2$ \\
2 & $0.2$ & $-5.6$ & $0.2-0.1\cdot(-1)=0.3$
\end{tabular}

Objective values after each step:
\[
f(0)=9,\qquad f(0.1)=8.41,\qquad f(0.2)=8.04,\qquad f(0.3)=7.69.
\]

The iterates move towards the optimum $w^{\star}=3$ using only the sign information.

\newpage
\section*{Assumptions}
\begin{assumption}[Smoothness]\label{ass:smooth}
$f$ is $L$–smooth, i.e.
\[
\forall x,y\in\R^{d}:\quad
f(y)\le f(x)+\nabla f(x)^{\top}(y-x)+\frac{L}{2}\|y-x\|_{2}^{2}.
\]
\end{assumption}
\begin{assumption}[Unbiased Median Sign]\label{ass:unbiased}
For every iteration $k$,
\[
\E\!\bigl[m_{k}\mid x_{k}\bigr]=\sign\!\bigl(\nabla f(x_{k})\bigr),
\]
where the expectation is taken over the random minibatch.
\end{assumption}
\begin{assumption}[Bounded Median Variance]\label{ass:var}
There exists $\sigma^{2}>0$ such that
\[
\E\!\bigl[\|m_{k}-\sign(\nabla f(x_{k}))\|_{2}^{2}\mid x_{k}\bigr]\;\le\;\sigma^{2},
\qquad\forall k.
\]
\end{assumption}
\begin{assumption}[Bounded Sign Vector]\label{ass:signbound}
For any $v\in\R^{d}$, $\|\sign(v)\|_{2}^{2}\le d$.
\end{assumption}
All constants $L,\sigma,d$ are known a priori; stepsize $\alpha$ will be chosen as a function of the horizon $T$.

\section*{Main Convergence Theorem}
\begin{theorem}\label{thm:main}
Assume \ref{ass:smooth}–\ref{ass:signbound} hold and let the stepsize be constant
\[
\alpha \;=\; \sqrt{\frac{2\bigl(f(x_{0})-f^{\star}\bigr)}{L d\,T}}\;,
\qquad f^{\star}:=\inf_{x}f(x).
\]
Then the iterates of SignSGD‑Median satisfy
\[
\frac{1}{T}\sum_{k=0}^{T-1}\E\!\bigl[\|\nabla f(x_{k})\|_{1}\bigr]
\;\le\;
\underbrace{\sqrt{\frac{2L d\bigl(f(x_{0})-f^{\star}\bigr)}{T}}_{\displaystyle O\!\left(T^{-1/2}\right)}\;+\;
\underbrace{\frac{L\alpha\sigma^{2}}{2}}_{\displaystyle O\!\left(T^{-1/2}\right)} .
\]
Consequently the average $\ell_{1}$–norm of the gradient converges to zero at the rate $O(T^{-1/2})$.
\end{theorem}

\section*{Theoretical Properties}
\begin{itemize}[leftmargin=*]
\item \textbf{Stability.} Because each update moves at most $\alpha$ in every coordinate (sign is bounded by $1$), the method is inherently robust to outlier gradients; the median further suppresses heavy‑tailed noise.
\item \textbf{Hyper‑parameter Sensitivity.} The bound depends on $\alpha\sqrt{d}$, showing that for high dimensional problems the stepsize should be scaled as $\alpha\propto 1/\sqrt{d}$, which matches empirical practice.
\item \textbf{Comparison to Baselines.}
    \begin{itemize}
        \item Compared to vanilla SGD (which attains $O(1/\sqrt{T})$ in $\ell_{2}$ norm), SignSGD‑Median obtains the same order but with communication cost $d$ bits per iteration.
        \item The median variance term $\sigma^{2}$ replaces the usual gradient variance $\sigma_{g}^{2}$; for heavy‑tailed noise the median yields a strictly smaller bound.
    \end{itemize}
\end{itemize}

\section*{Discussion}
The theorem shows that even though only one‑bit information per coordinate is transmitted, the algorithm retains the optimal $O(T^{-1/2})$ convergence rate for non‑convex smooth objectives, provided the median sign estimator is unbiased and has bounded variance. The dependence on the ambient dimension $d$ appears only through the factor $\sqrt{d}$, which is unavoidable when working with $\ell_{1}$‑type descent directions.

\newpage
\section*{Preliminaries (Definitions and Lemmas)}
\begin{lemma}[Smoothness Inequality]\label{lem:smooth}
If $f$ is $L$–smooth then for any $x,y\in\R^{d}$,
\[
f(y)\le f(x)+\nabla f(x)^{\top}(y-x)+\frac{L}{2}\|y-x\|_{2}^{2}.
\]
\end{lemma}
\begin{proof}
Standard result; follows from integrating the Lipschitz continuity of $\nabla f$. \qedhere
\end{proof}

\begin{lemma}[Sign Inner Product]\label{lem:signinner}
For any $v\in\R^{d}$,
\[
v^{\top}\sign(v)=\|v\|_{1}.
\]
\end{lemma}
\begin{proof}
Coordinate‑wise, $v_{j}\sign(v_{j})=|v_{j}|$. Summing over $j$ yields the claim. \qedhere
\end{proof}

\section*{Main Proof (Step‑by‑Step Derivation)}
We prove Theorem~\ref{thm:main}. Throughout the proof $\mathcal{F}_{k}$ denotes the $\sigma$‑algebra generated by $\{x_{0},\dots ,x_{k}\}$.

\bigskip
\noindent\textbf{[Step 1] Apply smoothness.}
Using Lemma~\ref{lem:smooth} with $x=x_{k}$ and $y=x_{k+1}=x_{k}-\alpha m_{k}$,
\[
f(x_{k+1})\le f(x_{k})-\alpha\,\nabla f(x_{k})^{\top}m_{k}
+\frac{L\alpha^{2}}{2}\|m_{k}\|_{2}^{2}. \tag{1}
\]

\noindent\textbf{[Step 2] Take conditional expectation.}
Condition on $\mathcal{F}_{k}$ and use Assumptions~\ref{ass:unbiased} and~\ref{ass:signbound}:
\[
\E\!\bigl[f(x_{k+1})\mid\mathcal{F}_{k}\bigr]
\le f(x_{k})-\alpha\,\E\!\bigl[\nabla f(x_{k})^{\top}m_{k}\mid\mathcal{F}_{k}\bigr]
+\frac{L\alpha^{2}}{2}\E\!\bigl[\|m_{k}\|_{2}^{2}\mid\mathcal{F}_{k}\bigr]. \tag{2}
\]

\noindent\textbf{[Step 3] Replace the inner product by an $\ell_{1}$ norm.}
From Assumption~\ref{ass:unbiased} and Lemma~\ref{lem:signinner},
\[
\E\!\bigl[\nabla f(x_{k})^{\top}m_{k}\mid\mathcal{F}_{k}\bigr]
= \nabla f(x_{k})^{\top}\sign\!\bigl(\nabla f(x_{k})\bigr)
= \|\nabla f(x_{k})\|_{1}. \tag{3}
\]

\noindent\textbf{[Step 4] Bound the second moment of $m_{k}$.}
Write $m_{k}= \sign(\nabla f(x_{k}))+ \Delta_{k}$ with $\Delta_{k}:=m_{k}-\sign(\nabla f(x_{k}))$.  
Using Assumption~\ref{ass:var} and the fact $\|\sign(\nabla f(x_{k}))\|_{2}^{2}\le d$ (Assumption~\ref{ass:signbound}),
\[
\E\!\bigl[\|m_{k}\|_{2}^{2}\mid\mathcal{F}_{k}\bigr]
= \|\sign(\nabla f(x_{k}))\|_{2}^{2}
   +\E\!\bigl[\|\Delta_{k}\|_{2}^{2}\mid\mathcal{F}_{k}\bigr]
\le d+\sigma^{2}. \tag{4}
\]

\noindent\textbf{[Step 5] Substitute (3) and (4) into (2).}
\[
\E\!\bigl[f(x_{k+1})\mid\mathcal{F}_{k}\bigr]
\le f(x_{k})-\alpha\|\nabla f(x_{k})\|_{1}
+\frac{L\alpha^{2}}{2}\bigl(d+\sigma^{2}\bigr). \tag{5}
\]

\noindent\textbf{[Step 6] Rearrange to isolate the gradient term.}
\[
\alpha\|\nabla f(x_{k})\|_{1}
\le f(x_{k})-\E\!\bigl[f(x_{k+1})\mid\mathcal{F}_{k}\bigr]
+\frac{L\alpha^{2}}{2}\bigl(d+\sigma^{2}\bigr). \tag{6}
\]

\noindent\textbf{[Step 7] Take full expectation and sum over $k=0,\dots ,T-1$.}
Denote $F_{k}:=\E[f(x_{k})]$. Summing (6) and using linearity of expectation yields
\[
\alpha\sum_{k=0}^{T-1}\E\!\bigl[\|\nabla f(x_{k})\|_{1}\bigr]
\le F_{0}-F_{T}
+ \frac{L\alpha^{2}T}{2}\bigl(d+\sigma^{2}\bigr). \tag{7}
\]
Since $F_{T}\ge f^{\star}$,
\[
\alpha\sum_{k=0}^{T-1}\E\!\bigl[\|\nabla f(x_{k})\|_{1}\bigr]
\le f(x_{0})-f^{\star}
+ \frac{L\alpha^{2}T}{2}\bigl(d+\sigma^{2}\bigr). \tag{8}
\]

\noindent\textbf{[Step 8] Divide by $\alpha T$ to obtain the average bound.}
\[
\frac{1}{T}\sum_{k=0}^{T-1}\E\!\bigl[\|\nabla f(x_{k})\|_{1}\bigr]
\le \frac{f(x_{0})-f^{\star}}{\alpha T}
+ \frac{L\alpha}{2}\bigl(d+\sigma^{2}\bigr). \tag{9}
\]

\noindent\textbf{[Step 9] Choose $\alpha$ to balance the two terms.}
Set
\[
\alpha \;=\; \sqrt{\frac{2\bigl(f(x_{0})-f^{\star}\bigr)}{L d\,T}},
\]
which satisfies $\alpha\le 1$ for sufficiently large $T$ (so that the bound on $\|m_{k}\|_{2}^{2}$ remains valid).  
Plugging this choice into (9) gives
\[
\frac{1}{T}\sum_{k=0}^{T-1}\E\!\bigl[\|\nabla f(x_{k})\|_{1}\bigr]
\le \sqrt{\frac{2L d\bigl(f(x_{0})-f^{\star}\bigr)}{T}}
   +\frac{L\alpha\sigma^{2}}{2}.
\]
The second term inherits the same $O(T^{-1/2})$ scaling because $\alpha=O(T^{-1/2})$.
This completes the proof of Theorem~\ref{thm:main}. \qedhere

\section*{Rate Derivation (With Constants)}
From the final inequality we can write explicitly
\[
\boxed{
\frac{1}{T}\sum_{k=0}^{T-1}\E\!\bigl[\|\nabla f(x_{k})\|_{1}\bigr]
\;\le\;
\underbrace{\sqrt{\frac{2L d\bigl(f(x_{0})-f^{\star}\bigr)}{T}}}_{\displaystyle C_{1}T^{-1/2}}
\;+\;
\underbrace{\frac{L\sigma^{2}}{2}\sqrt{\frac{2\bigl(f(x_{0})-f^{\star}\bigr)}{L d}}}_{\displaystyle C_{2}T^{-1/2}} }.
\]
Both $C_{1}$ and $C_{2}$ are problem‑dependent constants; the overall rate is $O(T^{-1/2})$.

\section*{Special Cases}
\begin{enumerate}[leftmargin=*]
\item \textbf{No variance ($\sigma^{2}=0$).} The bound reduces to the deterministic case
\[
\frac{1}{T}\sum_{k}\E[\|\nabla f(x_{k})\|_{1}]
\le \sqrt{\frac{2L d\bigl(f(x_{0})-f^{\star}\bigr)}{T}}.
\]
\item \textbf{High‑dimensional limit.} If $d$ grows while $L$ and the initial suboptimality remain fixed, the stepsize scales as $\alpha\sim d^{-1/2}$, guaranteeing that the per‑coordinate update magnitude stays $O(d^{-1/2})$.
\item \textbf{Convex $f$.} When $f$ is convex, $\|\nabla f(x_{k})\|_{1}$ lower‑bounds $f(x_{k})-f^{\star}$, yielding the familiar $O(T^{-1/2})$ convergence in function value.
\end{enumerate}

\end{document}